\documentclass[11pt]{article}
\usepackage[margin=2.4cm]{geometry}
\usepackage{amsmath,amssymb,graphicx,booktabs,xcolor}
\usepackage[colorlinks=true,linkcolor=black,urlcolor=blue]{hyperref}
% caption.sty unavailable in this TeX install; use the kernel default
\makeatletter\renewcommand{\@makecaption}[2]{\vskip 6pt
  \small\sbox\@tempboxa{\textbf{#1.} #2}%
  \ifdim \wd\@tempboxa >\hsize \textbf{#1.} #2\par
  \else \global\@minipagefalse \hb@xt@\hsize{\hfil\box\@tempboxa\hfil}\fi}\makeatother
\setlength{\parskip}{0.45em}\setlength{\parindent}{0pt}
\newcommand{\me}{\mathrm{me}}
\title{\vspace{-1.6cm}\textbf{The promoter lattice model}\\[2pt]
\large How the mathematics works, and why it cannot be bistable}
\author{\normalsize Notes for the GNP/MB H3K27me3 project}
\date{\normalsize 25 October 2026}

\begin{document}\maketitle\vspace{-0.9em}

\section{State space}

The lattice is $N$ sites (nucleosomes, 200\,bp pitch). Each site $i$ carries
exactly one state
\[
s_i \in \{\,\mathrm{U},\ \me 1,\ \me 2,\ \me 3,\ \mathrm{AC}\,\},
\]
so acetylation and methylation of H3K27 exclude one another by construction
rather than by a rule. H3K4me3 is a separate static field $k_i\in[0,1]$ taken
from the measured ChIP profile, not simulated.

\begin{figure}[h]\centering\includegraphics{fig_ladder.pdf}
\caption{States and transitions. Forward steps carry the relative rates
$8.7:1.0:1.3$ for $\mathrm{U}\!\to\!\me1\!\to\!\me2\!\to\!\me3$.}
\end{figure}

\section{Time stepping}

A fixed step $\Delta t = 30$\,s, not Gillespie. Each step applies discrete
enzyme events, then per-site field processes. $\Delta t$ is small enough that
every per-site probability stays well below one.

\section{Enzyme binding as a renewal process}

Each enzyme molecule alternates searching and bound, so one cycle costs
$\tau_{\mathrm{search}} + \tau_{\mathrm{bound}}$ and the number of binding
events in a step is
\[
N_{\mathrm{ev}} \sim \mathrm{Poisson}\!\left(\frac{n_{\mathrm{mol}}\,\Delta t}
{\tau_{\mathrm{search}}+\tau_{\mathrm{bound}}}\right).
\]
Each event selects a site from a weight vector $w$ ($\sum_i w_i=1$) and draws a
dwell $d\sim\mathrm{Exp}(\tau_{\mathrm{bound}})$. PRC2 has two binding modes:
\[
P(\text{stable}) = p_{\mathrm{stable}}\,(0.3 + 0.7\,\nu_i),
\qquad
d \sim \begin{cases}\mathrm{Exp}(10\,\mathrm{s}) & \text{stable}\\
\mathrm{Exp}(0.9\,\mathrm{s}) & \text{transient,}\end{cases}
\]
with $\nu_i$ the nucleation weight, so nucleation sites receive \emph{longer}
visits as well as more of them.

\section{Catalysis within a dwell}

Writing is first order over the dwell,
\[
P(\text{write}) = 1-e^{-k_i d},
\qquad
k_i = \frac{k_{\mathrm{cat}}}{60}\cdot
\frac{\rho_s}{\min_s \rho_s}\cdot \alpha_i \cdot \eta_i ,
\]
where $\rho = (8.7,\,1.0,\,1.3)$. Dividing by $\min_s\rho_s$ makes
$k_{\mathrm{cat}}$ refer to the \emph{rate-limiting} step; normalising to the
fastest step instead penalises the limiting one roughly ninefold, twice over.

\section{The two modulating fields}

Neighbours are coupled by a kernel that excludes self, decays over
$\lambda=2.5$ nucleosomes and is truncated at six, then row-normalised:
\[
K_{ij} \propto \exp\!\left(-\frac{|i-j|-1}{\lambda}\right),\quad
K_{ii}=0,\quad K_{ij}=0 \text{ for } |i-j|>6 .
\]

\begin{figure}[h]\centering\includegraphics{fig_kernel.pdf}
\caption{The coupling kernel. Read-write is strictly short range.}
\end{figure}

\textbf{Read-write} (allosteric stimulation of EED by neighbouring $\me3$):
\[
\alpha_i = 1 + (A-1)\,(Km)_i,\qquad m_j = \mathbf{1}[s_j=\me3].
\]
\textbf{Antagonism} by the active marks:
\[
\eta_i = \mathrm{clip}\big(1 - w_{\mathrm{inh}}\,(Ka)_i,\ 0,\ 1\big),\qquad
a_j = \mathrm{clip}\big(k_j + \mathbf{1}[s_j=\mathrm{AC}],\ 0,\ 1\big).
\]
Both multiply the catalytic rate. \textbf{Neither changes where PRC2 binds}, and
that turns out to decide everything in \S\ref{sec:mono}.

\section{Field processes and replication}

Per site per step, independent Bernoulli trials at rate $\times\,\Delta t$:
$\mathrm{AC}\to\mathrm{U}$ by HDAC; $\me3\to\me2\to\me1\to\mathrm{U}$ at the
basal rate $\delta$, iterated low to high so one random draw cannot cascade a
site down three levels in a single step; and exchange, which resets any state
to U at rate $k_{\mathrm{exch}}\,e_i$ with $e_i$ a Gaussian on the NDR (centre
$+150$\,bp, $\sigma=450$\,bp). Every 24\,h each site independently resets to U
with probability $0.5$.

\begin{figure}[h]\centering\includegraphics{fig_exch.pdf}
\caption{The exchange geometry matters more than it looks: at Ccnd2's me3 block
the weight is $2.4\times10^{-7}$ against $0.994$ at the NDR, so turnover acts
essentially only at the promoter core. The committed
$k_{\mathrm{exch}}=0.60/\mathrm{h}$ works \emph{because} it is NDR-local;
applied uniformly it erases every domain.}
\end{figure}

\section{The ensemble}

$C$ cells run independently, cell $c$ drawing
$n_{\mathrm{PRC2}}^{(c)}\sim\mathrm{lognormal}$ (CV $0.4$) and an active-zone
scale $\sim\mathcal{N}(1,0.25)$. Each cell's profile is time-averaged over the
second half of the run, then averaged across cells. That ensemble mean is what
is compared to bulk ChIP --- the comparison has to be made there, because bulk
ChIP cannot distinguish a uniform intermediate from a mixture.

\begin{figure}[h]\centering\includegraphics{fig_ensemble.pdf}
\caption{Ccnd2 with PRC2 seeded at the neighbouring promoter Gm56941. The
ensemble reproduces the measured block ($r=0.876$), and the cells agree with
each other everywhere except a ${\sim}1$\,kb zone at each domain edge.}
\end{figure}

\section{Why the model is monostable}\label{sec:mono}

Reduce to a single number $m$, the local $\me3$ fraction. Writing acts on
unmarked sites and is stimulated by neighbours; loss is first order:
\[
\frac{dm}{dt} \;=\; \underbrace{\beta\,(1-m)\big[1+(A-1)m\big]}_{W(m)}
\;-\; \underbrace{\gamma m}_{L(m)}
\;=\; -\beta(A\!-\!1)m^{2} + \big[\beta(A\!-\!2)-\gamma\big]m + \beta .
\]
This is a downward parabola with
\[
f(0) = +\beta > 0, \qquad f(1) = -\gamma < 0
\]
for \emph{any} positive $\beta,\gamma,A$. A continuous curve that is positive at
$0$ and negative at $1$ crosses zero exactly once on $(0,1)$. There is one fixed
point for every parameter choice: \textbf{the model is monostable, and no
allosteric gain can change that.}

Geometrically, $W(0)=\beta>0$ --- bare chromatin still gets methylated at the
basal rate --- and a straight loss line through the origin can cut such a curve
only once.

\begin{figure}[h]\centering\includegraphics{fig_flux.pdf}
\caption{Left: as built. Right: with recruitment. The only structural
difference is whether the writing curve passes through the origin.}
\end{figure}

This is confirmed empirically. Across 102 hysteresis pairs --- each parameter
set run from a fully marked domain and from a bare one --- the two starts always
settle at the same level.

\begin{figure}[h]\centering\includegraphics{fig_hysteresis.pdf}
\caption{Allosteric gain $1$--$50\times$, turnover $0.05$--$2.5/\mathrm{h}$ both
NDR-local and uniform, recruitment from strictly local to fully domain-wide,
replication on and off. Largest gap anywhere $0.043$, signs flipping at random.}
\end{figure}

\section{What would make it bistable}

Let $\me3$ \emph{recruit} PRC2 rather than only stimulate it. The writing flux
gains a factor $m$:
\[
\frac{dm}{dt} = m\,g(m), \qquad g(m) = \beta(1-m)\big[1+(A-1)m\big]-\gamma .
\]
Now $m=0$ is always a fixed point --- bare chromatin attracts no enzyme, so
nothing starts it --- and $g$ can vanish twice, giving three fixed points with
the outer two stable. The bistable window is
\[
1 \;<\; \frac{\gamma}{\beta} \;<\; 1 + \frac{(A-2)^{2}}{4(A-1)} ,
\]
the upper limit being the saddle-node fold where the two interior roots merge.

\begin{table}[h]\centering\small
\begin{tabular}{rrlll}
\toprule
$A$ & $\gamma/\beta$ & $m=0$ & interior roots & regime\\
\midrule
9  & 0.5 & unstable & $0.941$ & monostable high\\
9  & 1.5 & stable & $0.079,\ 0.797$ & \textbf{bistable}\\
9  & 2.0 & stable & $0.180,\ 0.695$ & \textbf{bistable}\\
9  & 4.0 & stable & --- & monostable off\\
25 & 2.0 & stable & $0.046,\ 0.913$ & \textbf{bistable}\\
25 & 4.0 & stable & $0.156,\ 0.803$ & \textbf{bistable}\\
25 & 8.0 & stable & --- & monostable off\\
\bottomrule
\end{tabular}
\caption{Fixed points of the recruitment-coupled model. Bistability requires
loss to outrun basal writing, and is a window rather than a guarantee.}
\end{table}

\begin{figure}[h]\centering\includegraphics{fig_bifurcation.pdf}
\caption{The fold sits at $\gamma/\beta = 2.53$ for $A=9$ and $6.51$ for
$A=25$: stronger read-write widens the window but does not remove the
requirement $\gamma>\beta$.}
\end{figure}

In the code the missing ingredient is one line. The binding weights
\texttt{w\_prc2} are computed once in \texttt{\_\_init\_\_} from nucleation and
never updated, so the enzyme lands in the same places regardless of what is
already marked. Closing that loop is a departure from what the in-cis
measurements support, so it should be framed as a hypothesis to test rather
than a correction.

\end{document}
